\centerline{\bf \Huge Домашнее задание}
\centerline{\bf \Huge Повороты}

\begin{flushright}
        \scriptsize{Несчастным поворот — на радость.\\ Уильям Шекспир. Король Лир}
\end{flushright}

\problem{1}
Внутри квадрата $A_1 A_2 A_3 A_4$ взята точка $P$. Из вершины $A_1$ опущен перпендикуляр на $A_2 P$, из $A_2$ - на $A_3 P$, из $A_3$ - на $A_4 P$, из $A_4$ - на $A_1 P$. Докажите, что все четыре перпендикуляра (или их продолжения) пересекаются в одной точке.

\problem{2}
На сторонах треугольника $A B C$ внешним образом построены правильные треугольники $A_1 B C, A B_1 C$ и $A B C_1$. Докажите, что $A A_1=B B_1=C C_1$.

\problem{3}
 На сторонах $C B$ и $C D$ квадрата $A B C D$ взяты точки $M$ и $K$ так, что периметртреугольника $C M K$ равен удвоенной стороне квадрата. Найдите величину угла $\angle MAK$.

\problem{4}
На сторонах $A B$ и $B C$ правильного треугольника $A B C$ взяты точки $M$ и $N$ так, что $M N \| A C,\  E$ --- середина отрезка $A N, D$ --- центр треугольника $B M N$. Найдите величины углов треугольника $C D E$.

\problem{5}
Точка $P$ лежит внутри равностороннего треугольника $A B C$. Докажите, что существует треугольник, стороны которого равны $P A, P B$ и $P C$.

\problem{6}
Вокруг квадрата описан параллелограмм. Докажите, что перпендикуляры, опущенные из вершин параллелограмма на стороны квадрата, образуют квадрат.

\problem{7}
Точки $P$ и $Q$ расположены внутри равностороннего треугольника $A B C$ так, что четырёхугольник $A P Q C$ - выпуклый, $A P=P Q=Q C$ и $\angle P B Q=30^{\circ}$. Докажите, что $A Q=B P$.

\problem{8}
\textbf{(Точка Торичелли.)} В треугольнике все углы не превосходят $120^{\circ}$.

a) Докажите, что внутри треугольника существует точка $T$, из которой все стороны видны под углами $120^{\circ}$.

б) Докажите, что точка $T$ является той точкой плоскости, для которой сумма расстояний от неё до вершин треугольника наименьшая.

в) Какой будет точка с наименьшей суммой расстояний до вершин треугольника для треугольника с углом, большим $120^{\circ}$.